\documentclass[12pt,a4paper]{article}
% Drake_preamble.tex - LaTeX Preamble Template
% 作者: 謝慕揚, MD, PhD, FESC
% 
% 使用說明:
% 1. 表格請使用 longtable，不要有垂直分隔線 |
% 2. 表格內換行使用 \newline，不要使用 \\
% 3. 藥物名稱、檢驗、檢查請維持英文
% 4. Reference 格式請使用 \href 加上驗證過的 DOI

% Including essential packages for Chinese typesetting
\usepackage{xeCJK}
\usepackage{fontspec}
% Configuring Chinese font with Noto Serif CJK TC, fallback to SimSun
\IfFontExistsTF{Noto Serif CJK TC}{
  \setCJKmainfont{Noto Serif CJK TC}
  \setCJKsansfont{Noto Serif CJK TC}
  \setCJKmonofont{Noto Serif CJK TC}
}{\setCJKmainfont{SimSun}
  \setCJKsansfont{SimSun}
  \setCJKmonofont{SimSun}
}

% Including other necessary packages
\usepackage{geometry}
\usepackage{titlesec}
\usepackage{enumitem}
\usepackage{xcolor}
\usepackage{colortbl}
\usepackage{tcolorbox}
\usepackage{booktabs}
\usepackage{longtable}
\usepackage{array}
\usepackage{graphicx}
\usepackage{tikz}
\usepackage{fancyhdr}
\usepackage{multicol}
\usepackage{datetime2}
\usepackage{hyperref}
\usepackage{mdframed}
\usepackage{amsmath}
\usepackage{amssymb}
\usepackage{multirow}

\hypersetup{
    colorlinks=true,
    linkcolor=emergency_blue,
    citecolor=emergency_blue,
    urlcolor=emergency_blue,
    pdfborder={0 0 0}
}

% Defining page geometry
\geometry{
  top=2.5cm,
  bottom=2.5cm,
  left=2cm,
  right=2cm,
  headheight=25pt
}

% Adjusting topmargin to compensate for increased headheight
\addtolength{\topmargin}{-10pt}

% Defining colors
\definecolor{emergency_red}{RGB}{186,24,27}
\definecolor{emergency_blue}{RGB}{0,114,188}
\definecolor{emergency_gray}{RGB}{120,120,120}
\definecolor{highlight_yellow}{RGB}{255,250,205}

% Setting title formats
\titleformat{\section}[block]{\Large\bfseries\color{emergency_red}}{\thesection}{10pt}{}
\titleformat{\subsection}[block]{\large\bfseries\color{emergency_blue}}{\thesubsection}{8pt}{}
\titleformat{\subsubsection}[block]{\normalsize\bfseries\color{emergency_gray}}{\thesubsubsection}{6pt}{}

% Defining custom environment for important notes
\newmdenv[
    linecolor=emergency_red,
    backgroundcolor=highlight_yellow,
    linewidth=2pt,
    topline=true,
    bottomline=true,
    leftline=true,
    rightline=true,
    innertopmargin=10pt,
    innerbottommargin=10pt,
    innerleftmargin=10pt,
    innerrightmargin=10pt
]{importantbox}

% Defining note command with tcolorbox
\newcommand{\note}[1]{
    \par
    \noindent
    \begin{tcolorbox}[
        colback=emergency_blue!5!white,
        colframe=emergency_blue,
        title=\textbf{重點提醒},
        fonttitle=\bfseries,
        boxrule=0.5pt,
        arc=3pt,
        left=6pt,
        right=6pt,
        top=6pt,
        bottom=6pt
    ]
    #1
    \end{tcolorbox}
}

% Key findings box
\newcommand{\keyfinding}[1]{
    \par
    \noindent
    \begin{tcolorbox}[
        colback=yellow!10!white,
        colframe=orange!75!black,
        title=\textbf{核心發現摘要},
        fonttitle=\bfseries,
        boxrule=0.5pt,
        arc=3pt
    ]
    #1
    \end{tcolorbox}
}

% Defining custom date format for ISO 8601
\DTMnewdatestyle{iso}{
  \renewcommand{\DTMdisplaydate}[4]{\number##1-\number##2-\number##3}
  \renewcommand{\DTMDisplaydate}[4]{\number##1-\number##2-\number##3}
}
\DTMsetdatestyle{iso}
\newcommand{\mydate}{\DTMtoday}


% Custom header for this document
\pagestyle{fancy}
\fancyhf{}
\fancyhead[L]{\textcolor{emergency_red}{NEJM Case Record 教學講義}}
\fancyhead[R]{\textcolor{emergency_blue}{產後腎病症候群：FSGS}}
\fancyfoot[C]{\textcolor{emergency_gray}{\thepage}}
\renewcommand{\headrulewidth}{1pt}
\renewcommand{\headrule}{\hbox to\headwidth{\color{emergency_red}\leaders\hrule height \headrulewidth\hfill}}

\begin{document}

\begin{center}
{\LARGE\bfseries\textcolor{emergency_red}{NEJM Case Records of the Massachusetts General Hospital}}\\[0.5cm]
{\Large\textbf{Case 1-2025: A 35-Year-Old Woman with Shortness of Breath and Edema in the Legs}}\\[0.3cm]
{\large 35歲女性，呼吸急促合併下肢水腫}\\[0.5cm]
{\normalsize N Engl J Med 2025;392:186-94. DOI: 10.1056/NEJMcpc2402498}\\[0.3cm]
{\normalsize 整理：謝慕揚 MD, PhD, FESC \quad 日期：\mydate}
\end{center}

\vspace{0.5cm}

%==============================================================================
\section{病例呈現 (Case Presentation)}
%==============================================================================

\subsection{主訴與現病史 (Chief Complaint \& History of Present Illness)}

A 35-year-old woman was admitted to this hospital because of shortness of breath and edema in the legs.

一位35歲女性因呼吸急促及下肢水腫入院。

\subsubsection{產前及生產經過 (Antepartum \& Delivery)}

Eight weeks before the current presentation, the patient was admitted for induction of labor at 39 weeks' gestation because of advanced maternal age and a fetus that was small for gestational age. She had begun taking low-dose aspirin at 13 weeks' gestation to decrease the risk of preeclampsia associated with antepartum obesity. She had no history of hypertension before or during pregnancy.

入院前八週，病患因高齡產婦及胎兒體重過輕，於妊娠39週時入院進行引產。她自妊娠13週起服用低劑量aspirin以降低與產前肥胖相關的子癇前症風險。病患於懷孕前及懷孕期間均無高血壓病史。

On admission for delivery:
\begin{itemize}[leftmargin=*]
    \item Blood pressure: 141/81 mm Hg（血壓：141/81 mm Hg）
    \item Pitting edema in legs present（下肢凹陷性水腫）
    \item Blood albumin: 2.1 g/dL (reference: 3.3--5.0)（血清白蛋白：2.1 g/dL，參考值：3.3--5.0）
    \item Urine protein/creatinine ratio: 4.1 (reference: <0.15)（尿蛋白/肌酸酐比值：4.1，參考值：<0.15）
    \item No headache, vision changes, or shortness of breath（無頭痛、視力改變或呼吸急促）
\end{itemize}

On the third hospital day, a healthy infant was born by uncomplicated vaginal delivery. Blood pressure after delivery was 122/88 mm Hg. She was discharged home on the fifth hospital day.

住院第三天，病患經由順利的陰道分娩產下一名健康嬰兒。產後血壓為122/88 mm Hg。於住院第五天出院返家。

\subsubsection{產後病程 (Postpartum Course)}

\begin{longtable}{>{\raggedright\arraybackslash}p{3cm} >{\raggedright\arraybackslash}p{12cm}}
\toprule
\textbf{Time Point} & \textbf{Clinical Events 臨床事件} \\
\midrule
\endfirsthead
\toprule
\textbf{Time Point} & \textbf{Clinical Events 臨床事件} \\
\midrule
\endhead
\midrule
\multicolumn{2}{r}{\textit{續下頁...}} \\
\endfoot
\bottomrule
\endlastfoot

After discharge & SBP <140 mm Hg at home; weight decreased by 9 kg in 2 weeks\newline 出院後居家收縮壓<140 mm Hg；2週內體重減輕9公斤 \\
\midrule

6 weeks before presentation & Headache developed, nearly constant, pain 3--10/10, relieved with ibuprofen\newline 出現頭痛，幾乎持續，疼痛程度3--10分，服用ibuprofen可緩解 \\
\midrule

4 weeks before presentation & Swelling in feet and legs with headache; weight increased by 6 kg over next week; SBP remained <140 mm Hg\newline 足部及下肢水腫合併頭痛；接下來一週體重增加6公斤；收縮壓維持<140 mm Hg \\
\midrule

1 week before presentation & Evaluated at another hospital ED for increased leg swelling\newline 因下肢水腫加劇至另一醫院急診就醫 \\
\midrule

6 days before presentation & Oral furosemide started; weight continued to increase; swelling worsened; dyspnea on exertion developed\newline 開始口服furosemide；體重持續增加；水腫惡化；出現運動性呼吸困難 \\

\end{longtable}

\subsection{過去病史與社會史 (Past Medical \& Social History)}

\begin{itemize}[leftmargin=*]
    \item \textbf{Obesity}: Pregravid BMI 40（肥胖：孕前BMI 40）
    \item \textbf{Anxiety}（焦慮症）
    \item \textbf{Allergy}: Penicillin (hives)（過敏：Penicillin導致蕁麻疹）
    \item \textbf{Social}: Married, lives with husband; no alcohol or illicit drugs; former occasional smoker, quit 1 year before presentation\newline 社會史：已婚，與丈夫同住；不飲酒、不使用毒品；曾偶爾吸菸，於就診前一年戒菸
    \item \textbf{Family history}: No family history of preeclampsia\newline 家族史：無子癇前症家族史
    \item \textbf{Prenatal screening} (9 months before): HIV-1/2 negative, HBV negative, HCV negative, HbA1c normal, thyroglobulin normal\newline 產前篩檢（9個月前）：HIV-1/2陰性、B型肝炎陰性、C型肝炎陰性、糖化血色素正常、甲狀腺球蛋白正常
\end{itemize}

\subsection{身體檢查 (Physical Examination)}

On current presentation:

\begin{longtable}{>{\raggedright\arraybackslash}p{4cm} >{\raggedright\arraybackslash}p{11cm}}
\toprule
\textbf{Vital Signs 生命徵象} & \textbf{Findings 發現} \\
\midrule
\endfirsthead
\toprule
\textbf{Vital Signs 生命徵象} & \textbf{Findings 發現} \\
\midrule
\endhead
\bottomrule
\endlastfoot

Temperature & 35.8°C \\
\midrule
Blood pressure & 142/85 mm Hg \\
\midrule
Pulse & 62 beats per minute \\
\midrule
Respiratory rate & 16 breaths per minute \\
\midrule
Oxygen saturation & 98\% on room air \\
\midrule
BMI & 45.2（體重較產後最低體重增加11公斤） \\
\midrule
Key finding & Pitting edema in the legs（下肢凹陷性水腫） \\

\end{longtable}

\subsection{檢驗數據 (Laboratory Data)}

\begin{longtable}{>{\raggedright\arraybackslash}p{4.5cm} >{\raggedright\arraybackslash}p{2.5cm} >{\raggedright\arraybackslash}p{2.5cm} >{\raggedright\arraybackslash}p{2.5cm} >{\raggedright\arraybackslash}p{3cm}}
\toprule
\textbf{Laboratory Test} & \textbf{8週前} & \textbf{1週前} & \textbf{目前} & \textbf{臨床意義} \\
\midrule
\endfirsthead
\toprule
\textbf{Laboratory Test} & \textbf{8週前} & \textbf{1週前} & \textbf{目前} & \textbf{臨床意義} \\
\midrule
\endhead
\midrule
\multicolumn{5}{r}{\textit{續下頁...}} \\
\endfoot
\bottomrule
\endlastfoot

Hemoglobin (12.0--16.0 g/dL) & 11.0 & 11.2 & 13.9 & 血液濃縮\newline Hemoconcentration \\
\midrule
Hematocrit (36.0--46.0\%) & 31.9 & 32.6 & 40.7 & 體液減少後上升 \\
\midrule
White cell count (4500--11000/μL) & 11,020 & 8,400 & 8,060 & 正常範圍 \\
\midrule
Platelet count (150,000--400,000/μL) & 262,000 & 297,000 & 272,000 & 正常，排除HELLP \\
\midrule
Creatinine (0.60--1.50 mg/dL) & 0.51 & 0.70 & 0.75 & 正常腎功能 \\
\midrule
BUN (8--25 mg/dL) & 10 & 22 & 22 & 輕度上升\newline 可能脫水 \\
\midrule
Total protein (6.0--8.3 g/dL) & 5.1 & 4.8 & 5.3 & 低蛋白血症 \\
\midrule
\textbf{Albumin (3.3--5.0 g/dL)} & \textbf{2.1} & \textbf{1.5} & \textbf{2.2} & \textbf{嚴重低白蛋白血症}\newline \textbf{Severe hypoalbuminemia} \\
\midrule
AST (9--32 U/L) & 22 & 24 & 32 & 正常 \\
\midrule
ALT (7--33 U/L) & 17 & 29 & 28 & 正常 \\
\midrule
TSH (0.40--5.00 μIU/mL) & -- & -- & 7.87 & 輕度上升\newline 需追蹤甲狀腺功能 \\
\midrule
Free T4 (0.9--1.8 ng/dL) & -- & -- & 0.9 & 正常低限 \\
\midrule
BNP & -- & Normal & Normal & 排除心衰竭 \\

\end{longtable}

\subsection{尿液檢查 (Urinalysis)}

\begin{longtable}{>{\raggedright\arraybackslash}p{5cm} >{\raggedright\arraybackslash}p{3cm} >{\raggedright\arraybackslash}p{3cm} >{\raggedright\arraybackslash}p{4cm}}
\toprule
\textbf{Test} & \textbf{Result} & \textbf{Reference} & \textbf{臨床意義} \\
\midrule
\endfirsthead
\bottomrule
\endlastfoot

Urine protein & 3+ & Negative & 顯著蛋白尿 \\
\midrule
Urine protein/creatinine ratio (8週前) & 4.1 & <0.15 & 腎病範圍蛋白尿\newline 相當於>4 g/day \\
\midrule
\textbf{Urine protein/creatinine ratio (目前)} & \textbf{5.2} & \textbf{<0.15} & \textbf{持續惡化的}\newline \textbf{腎病範圍蛋白尿} \\

\end{longtable}

\subsection{影像學檢查 (Imaging Studies)}

\begin{itemize}[leftmargin=*]
    \item \textbf{Chest radiography}: No abnormalities（胸部X光：無異常）
    \item \textbf{CT chest (PE protocol)}: Mild bronchial-wall thickening with scattered areas of mucous plugging; no pulmonary embolism（胸部電腦斷層（肺栓塞程序）：輕度支氣管壁增厚伴散在黏液栓塞；無肺栓塞）
    \item \textbf{Leg ultrasonography}: Negative for deep venous thrombosis（下肢超音波：無深部靜脈血栓）
\end{itemize}

%==============================================================================
\section{醫學英文寫作解析 (Medical English Writing Analysis)}
%==============================================================================

\begin{longtable}{>{\raggedright\arraybackslash}p{4.5cm} >{\raggedright\arraybackslash}p{3cm} >{\raggedright\arraybackslash}p{4cm} >{\raggedright\arraybackslash}p{3.5cm}}
\toprule
\textbf{原文例句} & \textbf{文法結構} & \textbf{中文翻譯} & \textbf{寫作技巧} \\
\midrule
\endfirsthead
\toprule
\textbf{原文例句} & \textbf{文法結構} & \textbf{中文翻譯} & \textbf{寫作技巧} \\
\midrule
\endhead
\midrule
\multicolumn{4}{r}{\textit{續下頁...}} \\
\endfoot
\bottomrule
\endlastfoot

A 35-year-old woman was admitted to this hospital because of shortness of breath and edema in the legs. &
Passive voice + because of + noun phrase &
一位35歲女性因呼吸急促及下肢水腫入院。 &
開場句型標準格式\newline 被動語態強調病患\newline 非入院者 \\
\midrule

Eight weeks before the current presentation, the patient was admitted for induction of labor. &
Time adverbial (fronted) + passive voice &
入院前八週，病患因引產而入院。 &
時間副詞置於句首\newline 強調時序關係 \\
\midrule

She had begun taking low-dose aspirin at 13 weeks' gestation to decrease the risk of preeclampsia. &
Past perfect + to-infinitive (purpose) &
她自妊娠13週起服用低劑量aspirin以降低子癇前症風險。 &
過去完成式表示\newline 更早發生的動作\newline to-infinitive表目的 \\
\midrule

The patient had no headache, vision changes, or shortness of breath. &
Subject + had no + noun phrase (series) &
病患無頭痛、視力改變或呼吸急促。 &
陰性發現的標準寫法\newline 使用or連接最後項目 \\
\midrule

On examination, the blood pressure was 142/85 mm Hg, and pitting edema in the legs was present. &
On + noun + passive construction &
檢查時，血壓為142/85 mm Hg，並可見下肢凹陷性水腫。 &
On examination開頭\newline 報告理學檢查發現 \\
\midrule

The blood albumin level was 2.2 g per deciliter (reference range, 3.3 to 5.0). &
Subject + be + value + (reference range) &
血清白蛋白為2.2 g/dL（參考範圍：3.3--5.0）。 &
檢驗值標準格式\newline 必附參考範圍 \\
\midrule

New-onset proteinuria in pregnancy is a diagnostic challenge and has implications for both maternal and fetal health. &
Noun phrase + be + noun + and + have + noun &
懷孕期間新發蛋白尿是診斷上的挑戰，對母體及胎兒健康均有影響。 &
討論段落開場\newline 點出臨床重要性 \\
\midrule

When preeclampsia is left untreated, it can progress to eclampsia and cause stroke, pulmonary edema, acute kidney injury, liver rupture, and seizures. &
When-clause + modal (can) + coordinate verbs &
若子癇前症未經治療，可進展至子癇並導致中風、肺水腫、急性腎損傷、肝臟破裂及癲癇。 &
條件句說明疾病進程\newline 列舉併發症 \\
\midrule

The correct diagnosis is crucial because management strategies for these conditions differ substantially. &
Subject + be + adjective + because-clause &
正確診斷至關重要，因為這些疾病的處置策略有顯著差異。 &
強調診斷重要性\newline because引出原因 \\
\midrule

Although primary membranous nephropathy is unlikely in this patient, another variant remains a possibility: lupus membranous nephropathy. &
Although-clause (concession) + colon + appositive &
雖然原發性膜性腎病變在此病患不太可能，但另一種變型仍有可能：狼瘡性膜性腎病變。 &
讓步句型+冒號\newline 引出具體診斷 \\
\midrule

A diagnosis of primary FSGS was favored given the diffuse podocyte foot-process effacement. &
Passive voice + given + noun phrase (reason) &
考量到廣泛的足細胞足突消失，傾向診斷原發性FSGS。 &
given引出支持診斷\newline 的證據 \\

\end{longtable}

%==============================================================================
\section{鑑別診斷 (Differential Diagnosis)}
%==============================================================================

\note{
\textbf{核心問題 (Core Clinical Question)}：\\
懷孕期間新發腎病範圍蛋白尿的病因為何？\\
What is the cause of new-onset nephrotic-range proteinuria in pregnancy?

\vspace{0.3cm}
\textbf{關鍵鑑別}：子癇前症（placental disease）vs. 腎絲球疾病（glomerular disease）
}

\subsection{妊娠期蛋白尿鑑別診斷架構}

\begin{importantbox}
\textbf{妊娠相關腎臟生理變化 (Pregnancy-Related Renal Changes)}：
\begin{itemize}[leftmargin=*]
    \item 低阻力胎盤循環發育 → 平均動脈壓下降、心輸出量及血容量增加
    \item 腎臟血漿流量增加 → 腎絲球過濾率(GFR)上升
    \item 即使GFR增加，懷孕期間蛋白尿仍屬異常
    \item 血清白蛋白在孕期會下降，但此病患第三孕期albumin 2.0 g/dL仍屬異常偏低
\end{itemize}
\end{importantbox}

\subsection{詳細鑑別診斷分析}

\begin{longtable}{>{\raggedright\arraybackslash}p{3.5cm} >{\raggedright\arraybackslash}p{4cm} >{\raggedright\arraybackslash}p{4cm} >{\raggedright\arraybackslash}p{3.5cm}}
\toprule
\textbf{鑑別診斷} & \textbf{支持證據} & \textbf{不支持證據} & \textbf{建議檢查} \\
\textbf{Differential Dx} & \textbf{Supporting} & \textbf{Against} & \textbf{Workup} \\
\midrule
\endfirsthead
\toprule
\textbf{鑑別診斷} & \textbf{支持證據} & \textbf{不支持證據} & \textbf{建議檢查} \\
\textbf{Differential Dx} & \textbf{Supporting} & \textbf{Against} & \textbf{Workup} \\
\midrule
\endhead
\midrule
\multicolumn{4}{r}{\textit{續下頁...}} \\
\endfoot
\bottomrule
\endlastfoot

\textbf{Preeclampsia}\newline 子癇前症 &
- SGA fetus（胎盤功能不全證據）\newline - 子癇前症危險因子（肥胖、高齡）\newline - 妊娠期最常見腎病範圍蛋白尿原因\newline - 入院時有一次BP 141/81 &
- 僅有一次輕度血壓升高\newline - 產後血壓正常\newline - 產後8週蛋白尿持續且惡化\newline - 典型子癇前症產後應迅速緩解 &
sFlt-1:PlGF ratio\newline （若早產且臨床不確定時有幫助） \\
\midrule

\textbf{Minimal change disease (MCD)}\newline 微小病變 &
- 成人腎病症候群原因之一\newline - 病患有使用NSAID（ibuprofen） &
- 蛋白尿早於NSAID使用\newline - 症狀為漸進性（數週），而非MCD典型的快速發作（數天）\newline - 年齡不典型（MCD多見於兒童） &
Kidney biopsy\newline 光學顯微鏡正常\newline 電子顯微鏡足突消失 \\
\midrule

\textbf{Primary membranous nephropathy}\newline 原發性膜性腎病變 &
- 成人最常見原發性腎病症候群之一\newline - 蛋白尿及病程符合 &
- 年齡不典型（多見於年長白人男性）\newline - 性別不符 &
Anti-PLA2R antibody\newline Kidney biopsy \\
\midrule

\textbf{Lupus membranous nephropathy}\newline 狼瘡性膜性腎病變（Class V lupus nephritis） &
- 好發於育齡女性\newline - 懷孕可誘發SLE發作或新發\newline - 高雌激素抑制細胞免疫、增加抗體產生 &
- 無腎外SLE表現\newline - 未提及ANA、anti-dsDNA、補體檢查\newline - Anti-PLA2R通常陰性 &
ANA, anti-dsDNA\newline C3, C4 complement\newline Kidney biopsy \\
\midrule

\textbf{Focal segmental glomerulosclerosis (FSGS)}\newline 局灶性節段性腎絲球硬化症 &
- 年齡及病程符合\newline - 成人特發性腎病症候群常見病因\newline - 急性發作的腎病症候群 &
- 需排除次發性FSGS（藥物、病毒感染） &
Kidney biopsy\newline （光學顯微鏡：節段性硬化\newline 電子顯微鏡：足突消失） \\
\midrule

\textbf{Gestational proteinuria}\newline 妊娠性蛋白尿 &
- 影響10\%孕婦 &
- 通常為非腎病範圍（<3.5 g/day）\newline - 本案蛋白尿程度過高\newline - 無高血壓 &
臨床追蹤\newline 30\%會進展為子癇前症 \\

\end{longtable}

\subsection{鑑別診斷推理過程 (Clinical Reasoning)}

\note{
\textbf{為何不是子癇前症？}

\begin{enumerate}[leftmargin=*]
    \item \textbf{血壓表現不典型}：子癇前症通常有更持續且嚴重的高血壓，此病患僅有一次輕度升高
    \item \textbf{產後病程}：子癇前症的蛋白尿應在產後6週內完全緩解，此病患產後8週仍有持續的腎病範圍蛋白尿
    \item \textbf{若未緩解}：通常是非腎病範圍且來自嚴重早產型子癇前症的殘餘腎臟損傷
\end{enumerate}
}

\begin{importantbox}
\textbf{sFlt-1:PlGF 比值的診斷價值}：

\begin{itemize}[leftmargin=*]
    \item 子癇前症源於血管新生失衡：病態胎盤釋放過量sFlt-1（可溶性VEGF受體）
    \item sFlt-1結合VEGF及PlGF → 廣泛內皮功能障礙
    \item \textbf{高sFlt-1:PlGF比值}：支持子癇前症
    \item \textbf{低sFlt-1:PlGF比值}：建議考慮其他腎絲球疾病
    \item 此生物標記已在歐洲及加拿大使用近十年，近期獲FDA核准於美國使用
\end{itemize}
\end{importantbox}

%==============================================================================
\section{診斷程序：腎臟切片 (Diagnostic Procedure: Kidney Biopsy)}
%==============================================================================

\subsection{光學顯微鏡發現 (Light Microscopy Findings)}

\begin{itemize}[leftmargin=*]
    \item 兩條皮質核心檢體，含24個腎絲球
    \item \textbf{4個腎絲球}顯示FSGS變化
    \item \textbf{1個腎絲球}含塌陷特徵（collapsing features）：
    \begin{itemize}
        \item 囊外上皮細胞增生（pseudocrescent）
        \item 圍繞疤痕及塌陷區域
        \item 蛋白重吸收滴（protein reabsorption droplets）
        \item 無壞死、纖維蛋白、發炎或腎絲球基底膜破裂
    \end{itemize}
    \item 反應性足細胞及部分內皮細胞腫脹
    \item PAS染色：腎絲球基底膜局灶性、節段性雙層化（duplication）→ 提示慢性內皮細胞損傷
    \item 無顯著間質發炎
    \item 近端腎小管上皮細胞顯示重度蛋白尿特徵：細胞質空泡化、蛋白重吸收滴
    \item 背景：輕度動脈硬化及小動脈玻璃樣變
\end{itemize}

\subsection{免疫螢光檢查 (Immunofluorescence)}

\begin{longtable}{>{\raggedright\arraybackslash}p{4cm} >{\raggedright\arraybackslash}p{4cm} >{\raggedright\arraybackslash}p{7cm}}
\toprule
\textbf{Marker} & \textbf{Result} & \textbf{臨床意義} \\
\midrule
\endfirsthead
\bottomrule
\endlastfoot

IgA, IgG, IgM & Negative & 排除免疫複合物沉積疾病 \\
\midrule
C3, C1q & Negative & 排除lupus nephritis \\
\midrule
Fibrin & Negative & 無血栓性微血管病變活性證據 \\
\midrule
Kappa, Lambda & Negative & 排除單株免疫球蛋白沉積 \\

\end{longtable}

\subsection{電子顯微鏡發現 (Electron Microscopy)}

\begin{itemize}[leftmargin=*]
    \item \textbf{廣泛且全面的足細胞足突消失}（diffuse and global podocyte foot-process effacement），幾乎涵蓋整個微血管表面
    \item 偶見腎絲球基底膜節段性雙層化
    \item 疤痕區域有少量電子緻密物及玻璃樣沉積
\end{itemize}

\note{
\textbf{廣泛足突消失的診斷意義}：\\
廣泛且全面的足細胞足突消失支持\textbf{原發性FSGS}的診斷。次發性FSGS（如因nephron loss導致）通常呈現局灶性足突消失。
}

%==============================================================================
\section{最終診斷 (Final Diagnosis)}
%==============================================================================

\begin{center}
\fbox{\parbox{0.9\textwidth}{
\centering
\Large\textbf{Postpartum Nephrotic Syndrome due to Focal Segmental Glomerulosclerosis}\\[0.3cm]
\large 產後腎病症候群，病因為局灶性節段性腎絲球硬化症
}}
\end{center}

\subsection{FSGS分類 (FSGS Classification)}

\begin{longtable}{>{\raggedright\arraybackslash}p{4cm} >{\raggedright\arraybackslash}p{5cm} >{\raggedright\arraybackslash}p{6cm}}
\toprule
\textbf{Category} & \textbf{Mechanism} & \textbf{此病患相關性} \\
\midrule
\endfirsthead
\bottomrule
\endlastfoot

\textbf{Primary FSGS} & 推測與通透性因子相關 & \textbf{最可能}：病程符合、廣泛足突消失、無明確次發原因 \\
\midrule
Secondary FSGS & Nephron loss後的病理模式 & 不符：通常呈現局灶性足突消失 \\
\midrule
Genetic FSGS & 遺傳性足細胞蛋白異常 & 可能性低：成人發病 \\
\midrule
FSGS of undetermined cause & 原因不明 & 需排除後考慮 \\

\end{longtable}

%==============================================================================
\section{治療計畫 (Treatment Plan)}
%==============================================================================

\subsection{初始治療 (Initial Treatment)}

\begin{longtable}{>{\raggedright\arraybackslash}p{3.5cm} >{\raggedright\arraybackslash}p{2cm} >{\raggedright\arraybackslash}p{1.5cm} >{\raggedright\arraybackslash}p{2cm} >{\raggedright\arraybackslash}p{5.5cm}}
\toprule
\textbf{Drug (Generic)} & \textbf{Dose} & \textbf{Route} & \textbf{Frequency} & \textbf{適應症/機轉} \\
\midrule
\endfirsthead
\toprule
\textbf{Drug (Generic)} & \textbf{Dose} & \textbf{Route} & \textbf{Frequency} & \textbf{適應症/機轉} \\
\midrule
\endhead
\midrule
\multicolumn{5}{r}{\textit{續下頁...}} \\
\endfoot
\bottomrule
\endlastfoot

Prednisone & High-dose & PO & Daily × 2 wks, then taper & 原發性FSGS第一線治療\newline 免疫抑制 \\
\midrule
Trimethoprim-sulfamethoxazole & Prophylactic dose & PO & Daily/3×wk & PJP預防\newline 高劑量類固醇使用期間 \\
\midrule
Omeprazole & 20--40 mg & PO & Daily & 腸胃道保護\newline 類固醇相關潰瘍預防 \\
\midrule
Calcium + Vitamin D & Standard & PO & Daily & 骨質保護\newline 類固醇誘發骨質疏鬆預防 \\

\end{longtable}

\subsection{治療反應與復發 (Treatment Response \& Relapse)}

\textbf{初始反應（治療4週後）}：
\begin{itemize}[leftmargin=*]
    \item 達到\textbf{部分緩解}（partial remission）
    \item 尿蛋白：6.2 g/day → 2.7 g/day
    \item 血清白蛋白：2.4 g/dL → 3.3 g/dL
\end{itemize}

\textbf{復發（Prednisone減量8週後）}：
\begin{itemize}[leftmargin=*]
    \item 下肢水腫增加
    \item 體重增加7 kg
    \item 尿蛋白：4.3 g/day
    \item Prednisone副作用：失眠、焦慮、疲倦、體重增加
\end{itemize}

\begin{importantbox}
\textbf{類固醇依賴型FSGS (Glucocorticoid-Dependent FSGS)}：\\
復發發生於類固醇減量過程中 → 定義為類固醇依賴型疾病\\
>50\%的部分緩解FSGS病患會復發
\end{importantbox}

\subsection{第二線治療 (Second-Line Treatment)}

\begin{longtable}{>{\raggedright\arraybackslash}p{3.5cm} >{\raggedright\arraybackslash}p{2cm} >{\raggedright\arraybackslash}p{1.5cm} >{\raggedright\arraybackslash}p{2cm} >{\raggedright\arraybackslash}p{5.5cm}}
\toprule
\textbf{Drug (Generic)} & \textbf{Dose} & \textbf{Route} & \textbf{Frequency} & \textbf{適應症/機轉} \\
\midrule
\endfirsthead
\bottomrule
\endlastfoot

Cyclophosphamide & Low-dose & PO & Daily × 2 months & 短期免疫抑制\newline 烷化劑 \\
\midrule
Rituximab & Standard & IV & 2 doses (間隔2週) then Q4M & Anti-CD20單株抗體\newline 持續B細胞耗竭 \\
\midrule
Prednisone & High-dose & PO & × 1 week, taper over 6 months & 類固醇依賴，需密切監測副作用 \\

\end{longtable}

\subsection{其他治療選項 (Alternative Options)}

根據KDIGO 2021指引，類固醇依賴型FSGS的其他選擇包括：
\begin{itemize}[leftmargin=*]
    \item Calcineurin inhibitors（Cyclosporine, Tacrolimus）
    \item Mycophenolate mofetil
    \item Cyclophosphamide
\end{itemize}

\subsection{所有FSGS病患的支持治療 (Supportive Care for All FSGS)}

\begin{itemize}[leftmargin=*]
    \item \textbf{RAAS blockade}：ACE inhibitor或ARB減少蛋白尿
    \item \textbf{血壓控制}
    \item \textbf{限鈉飲食}
\end{itemize}

\subsection{治療目標與預後 (Treatment Goals \& Outcome)}

\textbf{緩解定義}：
\begin{itemize}[leftmargin=*]
    \item \textbf{完全緩解}：尿蛋白<0.3 g/day + 穩定Cr + Albumin>3.5 g/dL
    \item \textbf{部分緩解}：蛋白尿顯著下降但未達完全緩解標準
\end{itemize}

\textbf{此病患預後}：
\begin{itemize}[leftmargin=*]
    \item 第二線治療6週後：再次達到部分緩解
    \item 尿蛋白降至1.4 g/day
    \item 血清白蛋白正常化
    \item 下肢水腫消退
    \item \textbf{最近追蹤}：完成2年rituximab治療，未再使用類固醇
    \item 維持部分緩解，尿蛋白1.2 g/day，腎功能正常
\end{itemize}

%==============================================================================
\section{學習重點 (Key Takeaways)}
%==============================================================================

\note{
\begin{enumerate}[leftmargin=*]
    \item \textbf{Clinical Pearl 1}: 懷孕期間新發蛋白尿需區分子癇前症（placental disease）與腎絲球疾病（glomerular disease），因處置策略截然不同。
    
    \textit{New-onset proteinuria in pregnancy requires differentiation between preeclampsia and glomerular diseases, as management differs substantially.}
    
    \item \textbf{Clinical Pearl 2}: 子癇前症的蛋白尿通常在產後6週內完全緩解。持續的腎病範圍蛋白尿應考慮其他腎絲球疾病。
    
    \textit{Proteinuria from preeclampsia typically resolves within 6 weeks postpartum. Persistent nephrotic-range proteinuria suggests an underlying glomerular disease.}
    
    \item \textbf{Clinical Pearl 3}: sFlt-1:PlGF比值可協助鑑別子癇前症與其他腎絲球疾病，尤其在臨床表現不典型且胎兒早產時特別有用。
    
    \textit{The sFlt-1:PlGF ratio can help distinguish preeclampsia from other glomerular diseases, particularly useful when clinical features are inconclusive and the fetus is preterm.}
    
    \item \textbf{Clinical Pearl 4}: 原發性FSGS以廣泛、全面的足細胞足突消失為特徵，腎臟切片是確診的關鍵。
    
    \textit{Primary FSGS is characterized by diffuse and global podocyte foot-process effacement; kidney biopsy is essential for diagnosis.}
    
    \item \textbf{Clinical Pearl 5}: 原發性FSGS合併腎病症候群的治療以免疫抑制為主，高劑量glucocorticoids為第一線。超過50\%達到部分緩解的病患會復發。
    
    \textit{Treatment of primary FSGS with nephrotic syndrome centers on immunosuppression, with high-dose glucocorticoids as first-line. More than 50\% of patients achieving partial remission will relapse.}
    
    \item \textbf{Clinical Pearl 6}: 類固醇依賴型FSGS可考慮rituximab，一種anti-CD20單株抗體，已顯示對此病患群的療效。
    
    \textit{For glucocorticoid-dependent FSGS, rituximab (anti-CD20 monoclonal antibody) has shown efficacy and can be considered as an alternative or adjunctive therapy.}
\end{enumerate}
}

\begin{importantbox}
\textbf{Red Flag（危險徵兆）}：

產後持續或惡化的腎病範圍蛋白尿 $\neq$ 子癇前症殘餘表現

\vspace{0.2cm}
\textbf{Persistent or worsening nephrotic-range proteinuria postpartum is NOT simply residual preeclampsia!}

\vspace{0.2cm}
必須考慮原發性腎絲球疾病並安排腎臟切片！
\end{importantbox}

\vspace{1cm}
\begin{center}
\textcolor{emergency_gray}{\rule{0.8\textwidth}{0.5pt}}\\[0.3cm]
\textit{Case presented at the Medicine Case Conference, Massachusetts General Hospital}\\
N Engl J Med 2025;392:186-94\\
整理：謝慕揚 MD, PhD, FESC
\end{center}

\end{document}